\begin{itemize}
    \item 盧卡斯定理 (Lucas's Theorem) :\\
      對於 $n, m \in \mathbb{Z}^{*}$ 與質數 $P$，
      $C(m,n) \pmod P$
        $= \prod ( C(m_i,n_i) )$
      其中 $m_i$ 為 $m$ 在 $P$ 進位下的第 $i$ 位數字。
    \item 斯特林近似公式 (Stirling approximation) : \\
      $n!\approx\sqrt{ 2 \pi n}(\frac{n}{e})^{n}e^\frac{1}{12n}$
    \item 第一類斯特林數 (Stirling Numbers of the first kind，大小為 $n$ 的置換中有 $k$ 個循環) : \\
      $S(n,k) = \text{多項式 } \prod_{i=0}^{n-1} (x+i) \text{ 中 } x^k \text{ 的係數}$
    \item 第二類斯特林數 (Stirling Numbers of the second kind，將 $n$ 個元素劃分為 $k$ 個非空集合) : \\
      $S(n,k) = \frac{1}{k!} \sum\limits_{j=0}^k (-1)^{k-j} {k \choose j} j^n$
    \item 皮克定理 (Pick's Theorem) : $A = i + b/2 - 1$\\
      $A$：多邊形面積、$i$：內部網格點數、$b$：邊上網格點數
    \item 卡塔蘭數 (Catalan number) : $C_n = {2n \choose n}/(n+1)$\\
      $C^{n+m}_{n}-C^{n+m}_{n+1} = (m+n)! \frac{n-m+1}{n+1}\quad \text{對於} \quad n \ge m$\\
      $C_n = \frac{1}{n+1}{2n \choose n} = \frac{(2n)!}{(n+1)!n!}$\\
      $C_0 = 1 \quad \text{且} \quad C_{n+1}= 2(\frac{2n+1}{n+2})C_n$\\
      $C_0 = 1 \quad \text{且} \quad C_{n+1} = \sum_{i=0}^{n} C_iC_{n-i} \quad \text{對於} \quad n \ge 0$
    \item 歐拉示性數 (Euler Characteristic): \\
      平面圖：$V-E+F-C=1$ \\
      凸多面體：$V-E+F=2$ \\
      $V,E,F,C$：分別為頂點數、邊數、面數（區域數）與連通分量數
    \item 基爾霍夫定理 / 矩陣樹定理 (Kirchhoff's theorem) : \\
      $A_{ii} = \deg(i), A_{ij} = (i,j) \in E\ ? -1 : 0$，
      刪除任意一行與一列後計算矩陣 $A$ 的行列式 $\det(A)$
    \item 波利亚計數定理 (Polya's theorem) ($c$ 為顏色數，$m$ 為置換中循環大小的數量): \\
      $\frac{1}{m} \sum_{i=1}^{m}{c^{\gcd(i,m)}}$
    \item 伯恩賽德引理 (Burnside's Lemma): \\
      $|X/G| = \frac{1}{|G|}\sum\limits_{g\in G} |X^g|$
    \item 錯排公式 (Derangements) ($n$ 個人中，每個人都不在原來位置的組合數): \\
      $dp[0]=1; dp[1]=0;$\\
      $dp[i]=(i-1)*(dp[i-1]+dp[i-2]);$
    \item 貝爾數 (Bell Number，將 $n$ 個人分組的方法總數) : \\
      $B_0= 1$\\
      $B_n= \sum_{k=0}^{n} S(n,k)\quad (\text{第二類斯特林數})$\\
      $B_{n+1}= \sum_{k=0}^{n}{n \choose k} B_k$
    \item 威爾遜定理 (Wilson's theorem) :\\
      $(p-1)! \equiv -1 \pmod p$
    \item 費馬小定理 (Fermat's little theorem) :\\
      $a^p \equiv a \pmod p$
    \item 歐拉降冪公式應用 (Euler's totient function):\\
      $A ^ {B ^ C} \pmod p = \text{pow}(A,\text{pow}(B,C,p-1)) \pmod p$
    \item 歐拉函數降冪公式 (Extended Euler's Theorem):\\
      $A^B \pmod C = A^{(B \bmod \phi(C)) + \phi(C)} \pmod C$
    \item 6 的倍數恆等式: \\
      $(a-1)^3 + (a+1)^3 + (-a)^3 + (-a)^3 = 6a$
    \item 標準楊表 (Standard Young Tableaux, SYT): \\
      $\lambda=(\lambda _{1}\geq \cdots \geq \lambda _{k})$，$\sum \lambda _{i} = n$ 記作 $\lambda \vdash n$\\
      $\lambda \vdash n$ 意指 $\lambda$ 為 $n$ 的整數分拆（Integer Partition），例如 $n = 10, \lambda = (6,4)$ 此分拆可表示一種楊表形狀。\\
      楊表：第 $1$ 列 $\lambda _{1}$ 格，$\cdots$，第 $k$ 列 $\lambda_{k}$ 格的方格圖。\\
      標準楊表：每列從左到右遞增，每行從上到下遞增。 \\
      若 $T$ 為某一排列（Permutation）經 RSK 對應後的標準楊表，則此排列的最長下降子序列（LDS）與最長遞增子序列（LIS）長度分別對應到 $T$ 的列數與行數。
    \item RSK 對應 (RSK Correspondence): \\
      一個排列與一對標準楊表 $(P,Q)$ 具有雙射關係。\\
      $P$：排列跑 RSK 演算法的結果，通常為半標準楊表（SSYT）。 \\
      $Q$：可用來還原排列（類似排列矩陣）。
    \item 鉤長公式 (Hook length formula，形狀為 $\lambda$ 的標準楊表個數): \\
      $f^{\lambda }={\frac {n!}{\prod h_{\lambda }(i,j)}}$ \\
      $h_{\lambda }(i,j)=$ 滿足 $(x = i \lor y = j) \land (x,y) \ge (i,j)$ 且 $(x,y)$ 落在形狀 $\lambda$ 內的點對 $(x,y)$ 數量。\\
      遞迴式：\\
      (i) $f^{(0,\cdots,0)} = 1$ \\
      (ii) $f^{(\lambda_{1},\cdots,\lambda_{m})} = \sum_{k=1}^{m} f^{(\lambda_{1},\cdots,\lambda_{k-1},\lambda_{k} - 1,\lambda_{k+1},\cdots,\lambda_{m})} $
    \item 環狀相鄰不同色著色數 (Chromatic polynomial of a cycle graph): \\
      $f(n,k) = (k-1)^n + (-1)^n(k-1)$\\
      其中 $n$ 為頂點數，$k$ 為顏色數。
\end{itemize}